\documentclass[11pt]{article}
\usepackage{xspace,mathrsfs}
\usepackage{graphicx,rotating}
\usepackage{url}
\usepackage{amscd,amsfonts,amsmath,amssymb}
\usepackage{color}
\usepackage{longtable}
\usepackage{multirow,setspace}
\renewcommand{\topfraction}{0.9}
\renewcommand{\textfraction}{0.1}
\usepackage{float}
\usepackage{booktabs}
\usepackage{caption}
\usepackage{fancyhdr}

\def\shownotes{0}   % set 1 for version with author notes
                    % set 0 for no notes

%%%%%%%%%%%%%%%%%%%%%%%%%%%%%%%%%%%%%%%%%%%%%%%%%%%%%%%%%%%%%%
%Some macros (you can ignore everything until "end of macros")
\topmargin 0pt \advance \topmargin by -\headheight \advance
\topmargin by -\headsep

\textheight 8.9in
\oddsidemargin 0pt \evensidemargin \oddsidemargin \marginparwidth
0.5in

\textwidth 6.5in

%%%%%%
\DeclareMathAlphabet{\mathsl}{OT1}{cmr}{m}{sl}
\DeclareMathAlphabet{\mathsc}{OT1}{cmr}{m}{sc}
\DeclareMathAlphabet{\mathslbf}{OT1}{cmr}{bx}{sl}

\newcommand{\procfont}[1]{\mathsc{#1}}
\newcommand{\algfont}[1]{{\mathsf{#1}}}
\newcommand{\schemefont}[1]{{\mathsf{#1}}}

\newcommand{\gamefont}[1]{\mathrm{#1}}
\newcommand{\getsr}{{{\,\leftarrow{\hspace*{-3pt}\raisebox{.75pt}{$\scriptscriptstyle\$$}}\,}}}
\newcommand{\bits}{\{0,1\}}
\newcommand{\Adv}{\mathbf{Adv}}
\newcommand{\AdvGen}[3]{\mathbf{Adv}^{\mathrm{#1}}_{#2}(#3)}
\newcommand{\Prob}[1]{{\Pr\left[\,{#1}\,\right]}}
\newcommand{\main}{\procfont{Game}}
\newcommand{\gameName}[1]{\underline{$\main$ #1} \\[2pt]}
\newcommand{\return}{\textrm{Ret}~}

\newcommand{\gamesfontsize}{\small}
\renewcommand{\captionfont}{}
\newcommand{\indd}{\hspace*{12pt}}

\newcommand{\mpage}[2]{\begin{minipage}{#1\textwidth} #2 \end{minipage}}
\newcommand{\fpage}[2]{\framebox{\begin{minipage}[t]{#1\textwidth}\setstretch{1.2}\gamesfontsize  #2 \end{minipage}}}
\newcommand{\hfpages}[3]{\hfpagess{#1}{#1}{#2}{#3}}
\newcommand{\hfpagess}[4]{
        \begin{tabular}{c@{\hspace*{.5em}}c}
        \framebox{\begin{minipage}[t]{#1\textwidth}\setstretch{1.2}\gamesfontsize #3 \end{minipage}}
        &
        \framebox{\begin{minipage}[t]{#2\textwidth}\setstretch{1.2}\gamesfontsize #4 \end{minipage}}
        \end{tabular}
    }

\newcommand{\Initialize}{\procfont{Initialize}}
\newcommand{\Finalize}{\procfont{Finalize}}
\newcommand{\InitializeCode}{\underline{\textrm{proc} $\Initialize$}\\[2pt]}
\newcommand{\FinalizeCode}[1]{\underline{\textrm{proc} $\Finalize(#1)$}\\[2pt]}

\newcommand{\calK}{{\cal K}}
\newcommand{\calM}{{\cal M}}

%%%%%%%  Author Notes %%%%%%%
\ifnum\shownotes=1
\newcommand{\authnote}[2]{{ $\ll$\textsf{\footnotesize #2}$\gg$}}
\else
\newcommand{\authnote}[2]{}
\fi
%%%%%%%%%%%%%%%%%%%%%%%%%%%%%%%%%

% end of macros
%%%%%%%%%%%%%%%%%%%%%%%%%%%%%%%%%%%%%%%%%%%%%%%%%%%%%%%%%%%%%


\title{COMPSCI 466: Homework 5}
\date{}

\usepackage{lastpage}
\pagestyle{fancy}
\lhead{\footnotesize \parbox{11cm}{COMPSCI 466, UMass Amherst, Spring 2026} }
\rhead{\footnotesize Instructor:  Adam O'Neill \\ \url{adamo@cs.umass.edu}}
\renewcommand{\headheight}{24pt}

\begin{document}

\maketitle
\thispagestyle{fancy}

\noindent Due April 8, 11:59pm.


\paragraph*{Problem 1.} (100 points.)
Let $D$ be the set of all strings whose length is a positive multiple of~$128$.
Define keyed hash function $H \colon \bits^{128} \times D \to \bits^{128}$
as follows for all $K \in \bits^{128}$ and $M \in D$:

\begin{center}
\begin{minipage}{1.5in}
\begin{tabbing}
123\=123\=123\=123\=123\=\kill
\textbf{Algorithm} $H_K(M)$: \\
\> Parse $M$ as $M[1] \| M[2] \| \cdots \| M[m]$ \\
\> $C[0] \gets 0^{128}$ \\
\> For $i = 1$ to $m$ do: \\
\>\> $B[i] \gets \mathsf{AES}_K\!\bigl(C[i-1] \oplus M[i]\bigr)$ \\
\>\> $C[i] \gets \mathsf{AES}_K\!\bigl(B[i]   \oplus M[i]\bigr)$ \\
\> Return $C[m]$
\end{tabbing}
\end{minipage}
\end{center}

\noindent
Above, $M$ is parsed as $m$ blocks of 128 bits each, the initial
chaining value is $C[0] = 0^{128}$, and each round applies
$\mathsf{AES}_K$ twice: first to $C[i-1] \oplus M[i]$ to produce $B[i]$,
then to $B[i] \oplus M[i]$ to produce the next chaining value $C[i]$.
The hash output is the final chaining value $C[m]$.

\medskip

\noindent (Part A --- 50 points.)
Show that $H$ is \emph{not} collision-resistant.
State your adversary $A$ in pseudocode.
Then provide a complete justification: trace through the full computation
of $H_K$ on both messages to verify they produce the same output,
explain why the two messages are distinct.
Analyze $\mathbf{Adv}^{\mathrm{cr}}_H(A)$  and
state the resource usage of $A$.

\medskip

\noindent (Part B --- 50 points.)
We say that $H$ is \emph{second-preimage resistant (SPR)} with respect to
the following game $\mathrm{SPR}^A_H$.

\begin{center}
\fpage{0.52}{
\begin{tabbing}
123\=123\=\kill
\InitializeCode
\> $K \getsr \bits^{128}$ \\
\> $M_1 \getsr \bits^{128}$ \\
\> \return $(K,\, M_1)$ \\[6pt]
\FinalizeCode{M_2}
\> \return $(M_2 \in D) \;\wedge\; (M_2 \neq M_1) \;\wedge\; \bigl(H_K(M_2) = H_K(M_1)\bigr)$
\end{tabbing}
}
\end{center}

\noindent
The adversary $A$ receives the output of $\Initialize$ and passes
its response to $\Finalize$.
We define
$\mathbf{Adv}^{\mathrm{spr}}_H(A) = \Pr\bigl[\mathrm{SPR}^A_H \Rightarrow 1\bigr]$,
where the probability is over the coins of $\Initialize$ and
any randomness of $A$.
We say $H$ is SPR if $\mathbf{Adv}^{\mathrm{spr}}_H(A)$ is negligible for
all efficient adversaries $A$.

\medskip\noindent
Show that $H$ is \emph{not} SPR.
State your adversary $A$ in pseudocode.
Prove that $H_K(M_2) = H_K(M_1)$ for your choice of $M_2$, tracing
through the full computation of $H_K$ on both inputs.
Argue that $M_2 \neq M_1$ and that $M_2 \in D$.
Finally, analyze $\mathbf{Adv}^{\mathrm{spr}}_H(A)$ and state the resource usage of $A$.

\end{document}
